% Eigenständiger Prosabeweis ohne lokale nummerierte Deklarationen.
\providecommand{\MogStar}{\mathbin{\ast}}
\providecommand{\MogDiamond}{\mathbin{\diamond}}

\section*{Die Rekonstruktionsfrage}
Zu einer Halbgruppe \((S,\MogStar)\) gehört die Potenzhalbgruppe ihrer
\emph{nichtleeren} Teilmengen. Ihr Träger und ihre Multiplikation sind
\[
  \PowNE{S}:=\powerset(S)\setminus\{\varnothing\},
  \qquad
  X\MogStar_{\mathcal P}Y
    :=\{x\MogStar y\mid x\in X,\ y\in Y\}.
\]
Das Produkt ist nichtleer, weil beide Faktoren nichtleer sind.
Die Assoziativität folgt aus der Assoziativität in \(S\): Beide
Klammerungen eines dreifachen Mengenprodukts bestehen genau aus den
Elementen \((x\MogStar y)\MogStar z=x\MogStar(y\MogStar z)\).

Ein Isomorphismus \(f\colon S\to T\) der Grundhalbgruppen induziert
durch \(X\mapsto f[X]\) einen Isomorphismus ihrer Potenzhalbgruppen.
Gilt auch die Umkehrung? Kann man also aus der Isomorphie der
Potenzhalbgruppen auf die Isomorphie von \(S\) und \(T\) schließen?

Das Mogiljanskaja-Gegenbeispiel beantwortet diese Frage ohne zusätzliche
Voraussetzungen mit Nein. Es gibt unendliche Halbgruppen mit
\[
  (S,\MogStar)\not\cong(T,\MogDiamond),
  \qquad
  (\PowNE{S},\MogStar_{\mathcal P})
    \cong(\PowNE{T},\MogDiamond_{\mathcal P}).
\]
Dies ist die Aussage von
\FormulaRefAuto{MogiljanskajaPairCounterexample}[thm] in Band~28.

Die Konstruktion fügt zu \(S\) genau einen neuen Punkt \(k\) hinzu.
Dieser Punkt unterscheidet die Grundhalbgruppen algebraisch. Auf der
Ebene der Teilmengen lässt sich dagegen eine unendliche Familie von
Teilmengen so umordnen, dass sie alle neuen \(k\)-haltigen Teilmengen
aufnimmt. Entscheidend ist, dass dabei die für die Multiplikation
maßgeblichen Daten erhalten bleiben und jede als Mengenprodukt
auftretende Menge unverändert bleibt.

\section*{Konstruktion und Beweis}
\begin{proof}[Beweis]
\textbf{Die beiden Grundhalbgruppen.}
Wir verwenden die natürlichen Zahlen mit \(0\in\mathbb N\) und setzen
\[
  \begin{aligned}
    a_i&:=(0,i),& L&:=\{a_i\mid i\in\mathbb N\},\\
    b_{ij}&:=(1,(i,j)),&
    D&:=\{b_{ij}\mid i,j\in\mathbb N\},\\
    \mathbf0&:=(2,0),& k&:=(3,0).
  \end{aligned}
\]
Die ersten Koordinaten sorgen dafür, dass \(L,D,\{\mathbf0\}\) und
\(\{k\}\) paarweise disjunkt sind. Innerhalb von \(L\) und \(D\)
sind die Indizes eindeutig: \(a_i=a_j\) gilt genau für \(i=j\),
und \(b_{ij}=b_{mn}\) genau für \(i=m\) und \(j=n\).
Diese speziellen Familien und ihre Eigenschaften werden im ersten
Kapitel der zugehörigen \MogProofsLink{} aufgebaut;
siehe \FormulaRefAuto{MogiljanskajaAElementsDef}[def],
\FormulaRefAuto{MogiljanskajaBElementsDef}[def] und
\FormulaRefAuto{MogiljanskajaLayersDisjoint}[thm]. Der fett gesetzte
Punkt \(\mathbf0\) ist ein Element der zu konstruierenden Halbgruppen,
nicht die natürliche Zahl \(0\).

Nun definieren wir
\[
  S:=L\cup D\cup\{\mathbf0\},
  \qquad T:=S\cup\{k\}.
\]
Auf beiden Trägern gilt dieselbe Multiplikationsregel:
\[
  a_i\MogStar a_j=b_{ij},\qquad
  a_i\MogDiamond a_j=b_{ij};
\]
jedes andere Produkt ist \(\mathbf0\). Dabei bezeichnet
\(\MogStar\) die Operation auf \(S\),
\(\MogDiamond\) diejenige auf \(T\). Anders formuliert:
Nur wenn beide Faktoren in \(L\) liegen, entsteht ein von
\(\mathbf0\) verschiedenes Produkt. Die eindeutigen Indizes und die
disjunkten Fälle machen die Vorschriften wohldefiniert. Alle Werte
liegen in \(D\cup\{\mathbf0\}\), also in beiden Trägern.
Dies ist die Konstruktion aus
\FormulaRefAuto{MogiljanskajaPairDef}[def].

Jedes Produkt zweier Elemente liegt außerhalb von \(L\). Ein weiteres
Multiplizieren ergibt daher stets \(\mathbf0\), unabhängig davon,
welches Paar man zuerst multipliziert. Für \(x,y,z\in S\) gilt also
\[
  (x\MogStar y)\MogStar z=\mathbf0
    =x\MogStar(y\MogStar z),
\]
und dieselbe Rechnung gilt in \(T\). Beide Operationen sind somit
assoziativ. Beide Halbgruppen sind unendlich, weil
\(i\mapsto a_i\) eine Injektion von \(\mathbb N\) in
\(L\subseteq S\subseteq T\) ist.

\medskip\noindent
\textbf{Warum die Grundhalbgruppen nicht isomorph sind.}
In \(S\) ist jedes Element außerhalb von \(L\) ein Produkt:
\(b_{ij}=a_i\MogStar a_j\) und
\(\mathbf0=\mathbf0\MogStar\mathbf0\).
In \(T\) tritt \(k\) dagegen niemals als Produkt auf, denn auch dort
liegen alle Produkte in \(D\cup\{\mathbf0\}\).

Angenommen, es gäbe einen Isomorphismus
\(f\colon(S,\MogStar)\to(T,\MogDiamond)\).
Wegen seiner Surjektivität besitzt \(k\) ein Urbild \(s\in S\).
Dieses Urbild kann nicht außerhalb von \(L\) liegen: Andernfalls wäre
\(s=p\MogStar q\) für geeignete \(p,q\in S\), und die
Homomorphiegleichung ergäbe
\[
  k=f(s)=f(p\MogStar q)=f(p)\MogDiamond f(q).
\]
Damit wäre \(k\) entgegen der Produktregel doch ein Produkt.
Also liegt \(s\) in \(L\), und es gibt ein \(i\in\mathbb N\)
mit \(f(a_i)=k\). Nun folgt
\[
  \begin{aligned}
    f(b_{i0})
      &=f(a_i\MogStar a_0)
        =k\MogDiamond f(a_0)=\mathbf0,\\
    f(b_{i1})
      &=f(a_i\MogStar a_1)
        =k\MogDiamond f(a_1)=\mathbf0.
  \end{aligned}
\]
Doch \(b_{i0}\neq b_{i1}\), weil \(0\neq1\). Dies widerspricht
der Injektivität von \(f\). Die Grundhalbgruppen sind daher nicht
isomorph. Zusammen mit Assoziativität und Unendlichkeit ist damit
\FormulaRefAuto{MogiljanskajaBasePairProperties}[thm] bewiesen.

\medskip\noindent
\textbf{Welche Informationen ein Mengenprodukt enthält.}
Für beliebige Teilmengen \(X,Y\subseteq T\) setzen wir
\[
  \begin{aligned}
    R(X,Y)&:=\{b_{ij}\mid i,j\in\mathbb N,
                            \ a_i\in X,\ a_j\in Y\},\\
    \varepsilon(X,Y)&:=
      \begin{cases}
        \varnothing,&X\subseteq L\text{ und }Y\subseteq L,\\
        \{\mathbf0\},&\text{sonst}.
      \end{cases}
  \end{aligned}
\]
Das \emph{Rechteck} \(R(X,Y)\) besteht aus allen Paarungen eines in
\(X\) vorkommenden \(a\)-Index mit einem in \(Y\) vorkommenden
\(a\)-Index. Es kann leer sein. Die Produktregel liefert unmittelbar
\begin{equation}
  X\MogStar_{\mathcal P}Y
    =R(X,Y)\cup\varepsilon(X,Y)
    \qquad(X,Y\in\PowNE{S}),
  \tag{*}\label{mog-reading-product-S}
\end{equation}
\begin{equation}
  X\MogDiamond_{\mathcal P}Y
    =R(X,Y)\cup\varepsilon(X,Y)
    \qquad(X,Y\in\PowNE{T}).
  \tag{**}\label{mog-reading-product-T}
\end{equation}
Denn zwei Faktoren aus \(L\) erzeugen genau die Elemente des
Rechtecks. Liegt mindestens ein Faktor außerhalb von \(L\), entsteht
\(\mathbf0\). Wenn etwa \(X\not\subseteq L\), können wir ein
\(x\in X\setminus L\) und, wegen der Nichtleerheit, ein \(y\in Y\)
nehmen; ihr Produkt ist \(\mathbf0\). Entsprechend gilt dies für
\(Y\not\subseteq L\). Liegen hingegen beide Mengen ganz in \(L\),
entsteht kein Nullprodukt. Hier ist die Voraussetzung nichtleerer
Teilmengen wesentlich. Dies erklärt
\FormulaRefAuto{MogiljanskajaPowerProductFormula}[thm].

Zum Beispiel ist für \(X=\{a_0,a_1,b_{00}\}\) und
\(Y=\{a_0,a_1\}\)
\[
  X\MogStar_{\mathcal P}Y
    =\{b_{00},b_{01},b_{10},b_{11},\mathbf0\}.
\]
Die vier \(b\)-Elemente kommen aus den Paarungen der
\(a\)-Elemente; der zusätzliche Faktor \(b_{00}\in X\)
erzeugt den Nullanteil.

\medskip\noindent
\textbf{Eine Reserve außerhalb der Produktrechtecke.}
Wir wollen Teilmengen von \(D\) verändern, ohne irgendein
Produktrechteck zu verändern. Dazu markieren wir drei Punkte:
\[
  c_0:=b_{00},\qquad c_1:=b_{11},\qquad c_2:=b_{10},
  \qquad P:=\{c_0,c_1\}.
\]
Als frei veränderlichen Bereich wählen wir alle Zeilen mit erstem
Index mindestens \(2\):
\[
  U:=\{b_{n+2,j}\mid n,j\in\mathbb N\},
  \qquad V:=U\cup\{c_2\},\qquad D':=P\cup V.
\]
In \(b_{n+2,j}\) läuft \(n\) durch \(0,1,2,\ldots\), der erste
Index somit durch \(2,3,4,\ldots\). Der zweite Index \(j\) ist
beliebig. Die Menge \(U\) umfasst also die vollständigen Zeilen ab
Zeile~\(2\). Es gilt
\(P\cap U=\varnothing\), \(c_2\notin P\cup U\) und
\(D'\subseteq D\). Der Punkt \(b_{01}\) gehört nicht zu \(D'\).
Diese Wahl stimmt mit
\FormulaRefAuto{MogiljanskajaDPrimeDef}[def] überein.

Unsere Reserve ist die Familie
\[
  \mathcal C:=\{P\cup W\mid W\subseteq V\}
    =\{K\subseteq D'\mid P\subseteq K\}.
\]
Jede ihrer Mengen enthält die beiden gegenüberliegenden Ecken
\(b_{00}\) und \(b_{11}\), aber niemals \(b_{01}\).
Deshalb liegt kein Produktrechteck in \(\mathcal C\): Enthielte
\(R(X,Y)\) die beiden Punkte \(b_{00}\) und \(b_{11}\),
so müssten \(a_0,a_1\) sowohl in \(X\) als auch in \(Y\)
liegen. Dann läge auch \(b_{01}\) in \(R(X,Y)\).
Ein solches Rechteck kann keine Teilmenge von \(D'\) sein.
Somit gilt für alle \(X,Y\subseteq T\)
\[
  R(X,Y)\notin\mathcal C.
\]
Auch \(\varnothing\) liegt nicht in \(\mathcal C\), weil jedes
Mitglied der Reserve den nichtleeren Kern \(P\) enthält.

\MogReserveDiagram

\medskip\noindent
\textbf{Zwei konkrete Verschiebungen.}
Um mit der Reserve alle neuen \(k\)-haltigen Teilmengen aufzunehmen,
benötigen wir Bijektionen
\(\sigma\colon U\bij V\) und \(\theta\colon U\bij D\).
Sie lassen sich ausdrücklich angeben:
\[
  \begin{aligned}
    \sigma(b_{n+2,j})&:=
      \begin{cases}
        b_{n+1,0},&j=0,\\
        b_{n+2,j},&j\neq0,
      \end{cases}\\[3pt]
    \theta(b_{n+2,j})&:=b_{nj}
       \qquad(n,j\in\mathbb N).
  \end{aligned}
\]
Die Abbildung \(\sigma\) verschiebt nur die nullte Spalte um einen
Platz. Dort ist \(c_2=b_{10}\) das Bild von \(b_{20}\), und
jeder Punkt \(b_{n+2,0}\) hat das eindeutige Urbild
\(b_{n+3,0}\). Alle Punkte mit zweitem Index ungleich \(0\)
bleiben fest. Die beiden Bereiche können keine gemeinsamen Bilder
haben. Also ist \(\sigma\) bijektiv auf \(V\).
Die Abbildung \(\theta\) verschiebt alle Zeilen um zwei Plätze:
Jeder Punkt \(b_{nj}\in D\) hat genau das Urbild \(b_{n+2,j}\).
Auch \(\theta\) ist daher bijektiv.

\MogShiftDiagram

Dies sind die in den \MogProofsLink{} bewiesenen Bijektionen
\FormulaRefAuto{MogiljanskajaSigmaBijective}[thm] und
\FormulaRefAuto{MogiljanskajaThetaBijective}[thm]. Ihre elementweise
Beschreibung macht deutlich, dass hier keine beliebige Auswahl von
Urbildern benötigt wird.

\medskip\noindent
\textbf{Wie die Reserve die neuen Teilmengen aufnimmt.}
Gesucht ist eine Bijektion von \(\powerset(D)\) auf
\(\powerset(D\cup\{k\})\), die jedes Produktrechteck festlässt.
Die Zielmenge enthält neben allen bisherigen Teilmengen von \(D\)
noch die neuen Mengen, die \(k\) enthalten. Wir bezeichnen deren
Familie mit
\[
  \mathcal E:=\{\{k\}\cup W\mid W\subseteq D\}.
\]
Somit ist \(\powerset(D\cup\{k\})=\powerset(D)\cup\mathcal E\),
und die beiden Familien rechts sind disjunkt. Außerhalb der Reserve
\(\mathcal C\) wollen wir alle Mengen unverändert lassen.
Es genügt deshalb, die Reserve bijektiv auf
\(\mathcal C\cup\mathcal E\) abzubilden: Ihre Bilder müssen
sowohl die bisherige Reserve als auch sämtliche neuen Mengen abdecken.

Dazu zerlegen wir \(\mathcal C\) in zwei Teilfamilien.
Die eine soll mit Hilfe von \(\sigma\colon U\bij V\) die gesamte
alte Reserve füllen; die andere soll mit Hilfe von
\(\theta\colon U\bij D\) alle neuen Mengen liefern.
Als Unterscheidungsmerkmal dient der Punkt \(c_2\):
\[
  \begin{aligned}
    \mathcal C_0&:=\{P\cup W\mid W\subseteq U\},\\
    \mathcal C_1&:=\{P\cup\{c_2\}\cup W\mid W\subseteq U\}.
  \end{aligned}
\]
Die Familie \(\mathcal C_0\) enthält genau die Reservemengen ohne
\(c_2\), die Familie \(\mathcal C_1\) genau diejenigen mit
\(c_2\). Wegen \(c_2\notin P\cup U\) sind sie disjunkt und
vereinigen sich zu \(\mathcal C\). In beiden Darstellungen ist
\(W\) eindeutig, nämlich der Schnitt der jeweiligen Menge mit \(U\).

Für \(W\subseteq U\) definieren wir nun
\[
  \begin{aligned}
    \psi(P\cup W)&:=P\cup\sigma[W],\\
    \psi(P\cup\{c_2\}\cup W)&:=\{k\}\cup\theta[W].
  \end{aligned}
\]
Die eindeutige Darstellung jeder Eingabemenge legt genau einen Wert
fest. Im ersten Zweig ist \(\sigma[W]\subseteq V\), also liegt
das Bild in \(\mathcal C\); im zweiten ist
\(\theta[W]\subseteq D\), also liegt es in \(\mathcal E\).
Damit ist \(\psi\) eine wohldefinierte Abbildung von
\(\mathcal C\) nach \(\mathcal C\cup\mathcal E\).

\MogAbsorptionDiagram

Wir zeigen nun: Jede Menge in \(\mathcal C\cup\mathcal E\)
besitzt genau ein Urbild in \(\mathcal C\).
Die Existenz eines Urbilds beweist die Surjektivität;
seine Eindeutigkeit beweist die Injektivität.
Da \(k\notin D\), sind \(\mathcal C\) und \(\mathcal E\)
disjunkt. Wir können die beiden Zielfamilien daher getrennt betrachten.

\emph{Der erste Zweig füllt die gesamte alte Reserve.}
Wir wählen eine beliebige Zielmenge \(H\in\mathcal C\).
Sie hat eindeutig die Form \(H=P\cup Z\), wobei
\(Z=H\setminus P\subseteq V\). Weil \(\sigma\colon U\bij V\)
bijektiv ist, gibt es genau eine Teilmenge \(W\subseteq U\)
mit \(\sigma[W]=Z\), nämlich \(W=\sigma^{-1}[Z]\).
Damit besitzt \(H\) genau das Urbild
\[
  K=P\cup\sigma^{-1}[Z]\in\mathcal C_0.
\]
Der erste Zweig erreicht also jedes Mitglied von \(\mathcal C\),
auch die Mengen, die \(c_2\) enthalten: Beispielsweise gilt
\[
  \psi(P)=P,\qquad
  \psi(P\cup\{b_{20}\})=P\cup\{c_2\}.
\]
Im zweiten Beispiel liegt das Urbild in \(\mathcal C_0\), sein Bild
aber in \(\mathcal C_1\). Das ist möglich, weil
\(\sigma(b_{20})=c_2\).

\emph{Der zweite Zweig füllt die neue Familie.}
Eine beliebige Zielmenge \(H\in\mathcal E\) hat eindeutig die Form
\(H=\{k\}\cup Z\) mit \(Z\subseteq D\).
Aus der Bijektivität von \(\theta\colon U\bij D\) folgt wieder
ein eindeutig bestimmter frei veränderlicher Anteil
\(W=\theta^{-1}[Z]\subseteq U\). Das einzige Urbild von \(H\) ist daher
\[
  K=P\cup\{c_2\}\cup\theta^{-1}[Z]\in\mathcal C_1.
\]
Zum Beispiel liefert der leere Anteil \(W=\varnothing\) die neue
Einermenge, während \(W=\{b_{20}\}\) noch \(b_{00}\) hinzufügt:
\[
  \psi(P\cup\{c_2\})=\{k\},\qquad
  \psi(P\cup\{c_2,b_{20}\})=\{k,b_{00}\}.
\]
Hier bezeichnet \(\sigma^{-1}[Z]\) beziehungsweise
\(\theta^{-1}[Z]\) jeweils das Bild von \(Z\) unter der
Umkehrfunktion der entsprechenden Bijektion.

Jede Zielmenge in \(\mathcal C\) besitzt somit genau ein Urbild
in \(\mathcal C_0\), und jede Zielmenge in \(\mathcal E\)
genau eines in \(\mathcal C_1\). Ein weiteres Urbild im jeweils
anderen Zweig ist ausgeschlossen, weil dessen Bild in der anderen,
disjunkten Zielfamilie liegt. Also hat jedes Element von
\(\mathcal C\cup\mathcal E\) genau ein Urbild in
\(\mathcal C_0\cup\mathcal C_1=\mathcal C\).
Damit ist die behauptete Bijektivität bewiesen:
\[
  \psi\colon\mathcal C\bij\mathcal C\cup\mathcal E.
\]
Eine Hälfte der Reserve füllt also die gesamte alte Reserve wieder
aus; die andere Hälfte liefert alle neuen Mengen. Dies ist die
konkrete Anwendung des Absorptionssatzes
\FormulaRefAuto{CoreFamilyAbsorptionBijective}[thm] aus Band~8.

Außerhalb der Reserve lassen wir jede Teilmenge von \(D\) fest:
\[
  \eta(K):=
    \begin{cases}
      K,&K\subseteq D,\ K\notin\mathcal C,\\
      \psi(K),&K\in\mathcal C.
    \end{cases}
\]
Jede Teilmenge von \(D\cup\{k\}\) gehört zu genau einer der
drei Familien
\[
  \powerset(D)\setminus\mathcal C,\qquad
  \mathcal C,\qquad\mathcal E.
\]
Die erste wird durch den Identitätszweig getroffen; die beiden anderen
werden durch die soeben konstruierte Bijektion \(\psi\) getroffen.
Damit ist
\[
  \eta\colon\powerset(D)\bij\powerset(D\cup\{k\}).
\]
Weil weder die leere Menge noch irgendein Produktrechteck in der
Reserve liegt, gilt insbesondere
\begin{equation}
  \eta(\varnothing)=\varnothing,
  \qquad \eta(R(X,Y))=R(X,Y).
  \tag{***}\label{mog-reading-fixed-rectangles}
\end{equation}
Außerdem ist wegen der Injektivität von \(\eta\)
\(\eta(K)=\varnothing\) genau dann erfüllt, wenn
\(K=\varnothing\). Damit sind die Eigenschaften aus
\FormulaRefAuto{MogiljanskajaReserveMapProperties}[thm] hergeleitet.

\medskip\noindent
\textbf{Die Bijektion zwischen den nichtleeren Potenzmengen.}
Wir halten die Schicht \(A_*:=L\cup\{\mathbf0\}\) fest und
wenden \(\eta\) nur auf den \(D\)-Anteil einer Menge an.
Für \(X\subseteq S\) sei
\[
  \widetilde\Phi(X):=(X\cap A_*)\cup\eta(X\cap D).
\]
Die disjunkten Zerlegungen
\(S=A_*\cup D\) und \(T=A_*\cup(D\cup\{k\})\)
zeigen, dass dies eine Bijektion von \(\powerset(S)\) auf
\(\powerset(T)\) ist. Eine ausdrückliche Umkehrfunktion ist
\[
  Y\longmapsto
    (Y\cap A_*)\cup\eta^{-1}\bigl(Y\cap(D\cup\{k\})\bigr).
\]
Bei beiden Zusammensetzungen erhält man auf \(A_*\) den unveränderten
Anteil und auf der anderen Schicht den ursprünglichen Anteil zurück.

Aus \(\eta(\varnothing)=\varnothing\) folgt
\(\widetilde\Phi(\varnothing)=\varnothing\).
Wegen der Injektivität ist die leere Menge auch das einzige Urbild
der leeren Menge. Daher bildet \(\widetilde\Phi\) jede nichtleere
Teilmenge von \(S\) auf eine nichtleere Teilmenge von \(T\) ab.
Umgekehrt besitzt jede nichtleere Teilmenge \(Y\) von \(T\)
ein Urbild; dieses kann nicht leer sein, da die leere Menge auf die
leere Menge abgebildet wird. Die Einschränkung von
\(\widetilde\Phi\) auf die nichtleeren Teilmengen, mit Zielmenge
\(\PowNE{T}\), ist folglich ebenfalls bijektiv. Wir nennen sie
\(\Phi\):
\[
  \Phi\colon\PowNE{S}\bij\PowNE{T},
  \qquad \Phi(X)=(X\cap A_*)\cup\eta(X\cap D).
\]
Das ist die Abbildung aus \FormulaRefAuto{MogiljanskajaPhiDef}[def];
der allgemeine Satz über diese Fortsetzung auf disjunkte Schichten ist
\FormulaRefAuto{LayerPowerMapNonemptyBijective}[thm] in Band~8.

\medskip\noindent
\textbf{Warum die Mengenprodukte erhalten bleiben.}
Die Abbildung \(\Phi\) lässt den \(L\)-Anteil und den
\(\mathbf0\)-Anteil jeder Menge unverändert. Ferner gilt
\[
  \Phi(X)\cap L=X\cap L,
  \qquad
  \Phi(X)\subseteq L\quad\Longleftrightarrow\quad X\subseteq L.
\]
Die erste Gleichheit folgt direkt aus der Definition von \(\Phi\):
Der Anteil in \(A_*\) bleibt erhalten, während
\(\eta(X\cap D)\subseteq D\cup\{k\}\) keinen Punkt von
\(L\) enthält. Ebenso gilt
\(\mathbf0\in\Phi(X)\) genau dann, wenn \(\mathbf0\in X\).

Für die zweite Aussage betrachten wir die Punkte außerhalb von
\(L\). In \(S\) sind das \(\mathbf0\) und die Punkte von
\(D\); in \(T\) kommt noch \(k\) hinzu. Der entsprechende
\(D\cup\{k\}\)-Anteil von \(\Phi(X)\) ist genau
\(\eta(X\cap D)\). Deshalb gilt Schritt für Schritt
\[
  \begin{aligned}
    \Phi(X)\subseteq L
      &\ \Longleftrightarrow\
        \mathbf0\notin\Phi(X)\ \text{und}\
        \Phi(X)\cap(D\cup\{k\})=\varnothing\\
      &\ \Longleftrightarrow\
        \mathbf0\notin X\ \text{und}\
        \eta(X\cap D)=\varnothing\\
      &\ \Longleftrightarrow\
        \mathbf0\notin X\ \text{und}\ X\cap D=\varnothing\\
      &\ \Longleftrightarrow\ X\subseteq L.
  \end{aligned}
\]
Im dritten Schritt verwenden wir, dass \(\eta\) genau die leere
Menge auf die leere Menge abbildet.

Seien nun \(X,Y\in\PowNE{S}\). Für jeden Index \(i\) gilt
\(a_i\in\Phi(X)\) genau dann, wenn \(a_i\in X\);
entsprechend gilt für jeden Index \(j\), dass
\(a_j\in\Phi(Y)\) genau dann, wenn \(a_j\in Y\).
Die Definition des Rechtecks liefert daher
\[
  \begin{aligned}
    b_{ij}\in R(\Phi(X),\Phi(Y))
      &\ \Longleftrightarrow\ a_i\in\Phi(X)\ \text{und}\ a_j\in\Phi(Y)\\
      &\ \Longleftrightarrow\ a_i\in X\ \text{und}\ a_j\in Y\\
      &\ \Longleftrightarrow\ b_{ij}\in R(X,Y).
  \end{aligned}
\]
Beide Rechtecke sind Teilmengen von \(D\), also folgt
\[
  R(\Phi(X),\Phi(Y))=R(X,Y).
\]
Für die Nullanteile betrachten wir die beiden Fälle in der Definition
von \(\varepsilon\). Der Wert \(\varepsilon(\Phi(X),\Phi(Y))\) ist genau
dann leer, wenn \(\Phi(X)\subseteq L\) und
\(\Phi(Y)\subseteq L\) gelten. Nach der eben bewiesenen
Äquivalenz ist dies genau dann der Fall, wenn \(X\subseteq L\)
und \(Y\subseteq L\) gelten, also wenn
\(\varepsilon(X,Y)=\varnothing\) ist. In allen anderen Fällen
sind beide Nullanteile gleich \(\{\mathbf0\}\). Somit gilt
\[
  \varepsilon(\Phi(X),\Phi(Y))=\varepsilon(X,Y).
\]
Da \(\Phi(X)\) und \(\Phi(Y)\) nichtleer sind, dürfen wir die
Produktformel \eqref{mog-reading-product-T} auf sie anwenden.
Mit den beiden soeben bewiesenen Gleichheiten und anschließend
\eqref{mog-reading-product-S} erhalten wir
\begin{equation}
  \begin{aligned}
    \Phi(X)\MogDiamond_{\mathcal P}\Phi(Y)
      &\overset{\eqref{mog-reading-product-T}}{=}
        R(\Phi(X),\Phi(Y))\cup\varepsilon(\Phi(X),\Phi(Y))\\
      &=R(X,Y)\cup\varepsilon(X,Y)\\
      &\overset{\eqref{mog-reading-product-S}}{=}
        X\MogStar_{\mathcal P}Y.
  \end{aligned}
  \tag{IV}\label{mog-reading-product-images}
\end{equation}

Es bleibt zu prüfen, was \(\Phi\) mit dem Produkt selbst macht.
Setzen wir \(Z:=X\MogStar_{\mathcal P}Y\), so ist nach der
Produktformel
\[
  Z\cap D=R(X,Y),\qquad
  Z\cap A_*=\varepsilon(X,Y).
\]
Hier benutzen wir, dass \(R(X,Y)\subseteq D\), während der
Nullanteil eine Teilmenge von \(\{\mathbf0\}\) ist.
Da \(\eta\) jedes Produktrechteck festlässt, erhalten wir
\begin{equation}
  \begin{aligned}
    \Phi(Z)
      &=(Z\cap A_*)\cup\eta(Z\cap D)\\
      &=\varepsilon(X,Y)\cup\eta(R(X,Y))\\
      &\overset{\eqref{mog-reading-fixed-rectangles}}{=}
        \varepsilon(X,Y)\cup R(X,Y)=Z.
  \end{aligned}
  \tag{V}\label{mog-reading-fixed-product}
\end{equation}
Zusammen ergibt dies die Homomorphiegleichung
\[
  \Phi(X\MogStar_{\mathcal P}Y)
    \overset{\eqref{mog-reading-fixed-product}}{=}
      X\MogStar_{\mathcal P}Y
    \overset{\eqref{mog-reading-product-images}}{=}
      \Phi(X)\MogDiamond_{\mathcal P}\Phi(Y).
\]
Die Bijektion \(\Phi\) ist somit ein Isomorphismus der
Potenzhalbgruppen, wie in
\FormulaRefAuto{MogiljanskajaPowerSemigroupIsomorphism}[thm]
behauptet. Ihre Grundhalbgruppen sind nach dem ersten Teil des
Beweises unendlich und nicht isomorph. Damit ist das Gegenbeispiel
vollständig bewiesen.
\end{proof}

\section*{Was die Konstruktion sichtbar macht}
Die Grundhalbgruppen unterscheiden sich daran, wie ein einzelner Punkt
in Produkte eingeht: Das zusätzliche Element \(k\) ist selbst kein
Produkt, vernichtet aber jeden weiteren Faktor zu \(\mathbf0\).
In \(S\) gibt es kein Element, das unter einem Isomorphismus seine
Rolle übernehmen könnte.

Bei den Potenzhalbgruppen ist ein Element hingegen eine ganze
nichtleere Teilmenge. Ein Isomorphismus der Potenzhalbgruppen muss
die Anzahl der Elemente einer solchen Teilmenge nicht erhalten.
Unsere Abbildung liefert ein konkretes Beispiel:
\[
  \Phi(\{b_{00},b_{11},b_{10}\})=\{k\}.
\]
Denn die linke Menge ist \(P\cup\{c_2\}\), also ein Mitglied
von \(\mathcal C_1\) mit leerem frei veränderlichem Anteil;
der zweite Zweig von \(\psi\) schickt sie auf \(\{k\}\).
Zugleich fixiert \(\Phi\) jede Einermenge \(\{s\}\) mit
\(s\in S\): Punkte aus \(A_*\) bleiben ohnehin fest, und eine
Einermenge in \(D\) gehört nicht zur Reserve, weil sie nicht beide
Punkte von \(P\) enthält. Die neue Einermenge \(\{k\}\) hat
also eine dreielementige Menge als Urbild.
Diese dreielementige Menge ist kein Produktrechteck, da die Ecke
\(b_{01}\) fehlt. Die Multiplikation lässt somit gerade dort
genügend Spielraum für die Umordnung, wo sie keine Rechtecke erzeugt.

Die Konstruktion benötigt nur die ausdrücklich angegebenen
Verschiebungen und ihre Umkehrfunktionen. Ihr Kern ist nicht eine
beliebige Bijektion zwischen zwei Potenzmengen, sondern eine solche
Bijektion, die alle Produktrechtecke und die leere Menge festlässt.
Erst diese zusätzlichen Eigenschaften machen daraus den gesuchten
Isomorphismus der Potenzhalbgruppen.

Das Gegenbeispiel betrifft unendliche Halbgruppen; die entsprechende
Rekonstruktionsfrage für endliche Halbgruppen wird dadurch nicht entschieden.
Die Konstruktion, ihre Hilfsaussagen und sämtliche tabellarischen
Ableitungen stehen in den zugehörigen \MogProofsLink{}.
Dorthin führen auch die Verweise mit dem Nummernpräfix \(28\mathrm E\).

\section*{Originalarbeit und vorliegende Fassung}
Das Gegenbeispiel geht auf E.~M.~Mogiljanskaja zurück:
\href{https://doi.org/10.1007/BF02389140}
{\emph{Non-isomorphic Semigroups with Isomorphic Semigroups of Subsets}},
\emph{Semigroup Forum}~6 (1973), S.~330--333.

Die Originalarbeit behandelt freie nilpotente Halbgruppen der Klasse
\(n\geq3\). Unsere Grundhalbgruppe entspricht dem Fall \(n=3\):
Die von \(\mathbf0\) verschiedenen Elemente entsprechen Wörtern
der Länge eins oder zwei, und jedes Produkt aus drei Faktoren ist
\(\mathbf0\). Hinzu kommt das neue Element \(k\).

Der Unterschied liegt in der Umordnung der Teilmengen.
Mogiljanskaja verwendet eine Reservefamilie aus Teilmengen geeignet
gewählter Wörter der Länge \(n-1\). Auf S.~331 begründet sie die
Existenz der benötigten Bijektion durch die gleiche Mächtigkeit der
beteiligten Familien und wählt eine solche Bijektion, ohne eine
konkrete Formel anzugeben.

Hier verwenden wir eine andere Reservefamilie: Ihre Mengen enthalten
die beiden festen Marker \(b_{00}\) und \(b_{11}\) und lassen
\(b_{01}\) aus. Die Verschiebungen \(\sigma\) und \(\theta\)
liefern dafür die ausdrücklich angegebene Abbildung \(\psi\).
Die vorliegende Fassung führt somit den Grundgedanken des
Originalbeweises in einem Spezialfall aus, verwendet für die
Reservefamilie und ihre Bijektion aber eine eigene konkrete
Konstruktion.
